\scanpage{79}
% printed form
% letterhead
UNIVERSITE\hfill Montpellier, le 8 Décembre 1977\\
DES SCIENCES ET TECHNIQUES\\
DU LANGUEDOC\\
------------\\
Service de la Scolarité\\
JP/GJ\\
Diffusion : Liste A\\
N°7041

\uline{AVIS DE SOUTENANCE D'UNE THESE DE DOCTORAT DE}

\uline{SPECIALITE : ELECTRONIQUE DES SOLIDES}

\uline{Nom et prénom du candidat} : DAROLLES Jean-Marie

\uline{Titre de la thèse} : Etude de la conversion de l'énergie à partir de thermoéléments réfractaires Bore Silicium.

\uline{Date, heure et lieu de soutenance prévus} :\\
12 DECEMBRE 1977 à 9 Heures\\
Salle de réunion du C.E.E.S., USTL

\uline{Nom des directeurs de thèse} :\\
J.L ROBERT et J.M. DUSSEAU

\uline{Désignation du laboratoire où la thèse est préparée} :\\
Laboratoire d'Electronique C.E.E.S

\uline{Analyse de l'ouvrage} : Après avoir défini les grandeurs thermoélectriques classiques, nous montrons comment on peut optimiser le rendement d'un convertisseur. Des montages permettant d'effectuer des mesures de 77 à 1600 K sont décrits. Les résultats concernant les composés réfractaires Bore-Silicium montrent l'intérêt d'utiliser de tels corps dans la conversion thermoélectrique.

\uline{Durée approximative des recherches ayant abouti à la thèse} : 2 ans

\uline{Composition du jury proposé} :\\
MM. : B. PISTOULET, Président\\
L. LASSABATERE\\
J.-L. ROBERT\\
C. COMBESCURE\\
J.-M. DUSSEAU

Le Secrétaire Général,\\
\handwritten{\unsure{E. Siau|E. Si}}\\
\uline{E. SIAU}

% archive sheet number: 78
